\section{Related Work}

\subsection{Generative Simulation for Robotics}
Simulation is a cornerstone for scaling robot learning, enabling safe and rapid data collection~\cite{todorov2012mujoco, makoviychuk2021isaac}. Early approaches relied on procedural generation with fixed asset libraries to create diverse environments~\cite{deitke2023procthor, szot2021habitat}, but these methods are inherently limited by the diversity of hand-crafted asset databases and often lack photorealistic textures. Recent works have begun to leverage generative models to automate simulation design at a higher level. GenSim~\cite{wang2024gensim} and RoboGen~\cite{wang2023robogen} utilize LLMs to generate task codes and scene configurations, Holodeck~\cite{yang2024holodeck} compositionally arranges assets into 3D environments from language descriptions, and MimicGen~\cite{mandlekar2023mimicgen} automatically scales demonstrations to new object placements. However, all these methods are fundamentally constrained by fixed asset repositories, limiting visual variety and semantic open-endedness. V-Dreamer overcomes this by synthesizing novel 3D geometries and photorealistic textures entirely from scratch via diffusion models, enabling truly open-vocabulary scene diversity unconstrained by any asset library.

\subsection{Video Generation as Motion Priors}
The rapid progress in video generation~\cite{blattmann2023stable, wu2023tune} has sparked interest in using video models as world simulators or motion priors for robotics. UniSim~\cite{yang2023unisim} learns a universal interactive simulator from real-world action-annotated videos. Video language planning~\cite{cohen2024avdc} and UniPi~\cite{du2023learning} use video diffusion models to imagine future frames as a planning signal. GR-2~\cite{gu2024gr2} pre-trains on large-scale internet video and fine-tunes for manipulation, showing that video prediction priors enhance generalization. Genie~\cite{bruce2024genie} learns controllable world models from unlabeled play videos as interactive training grounds. Despite this progress, most approaches operate in 2D pixel space, making it difficult to extract precise, physically grounded 3D trajectories for robot control. V-Dreamer addresses this by employing a Sim-to-Gen alignment module based on CoTracker3~\cite{karaev2023cotracker} and VGGT~\cite{wang2025vggt} to lift 2D visual motion into physically executable end-effector trajectories, directly bridging visual synthesis and kinematic control.


\subsection{Sim-to-Real Transfer and Generalization}
Bridging the reality gap between simulation and the physical world remains a central challenge in robot learning. Domain Randomization (DR)~\cite{tobin2017domain, peng2018sim} randomizes physical and visual parameters to expose policies to diverse conditions, but requires manual tuning of randomization ranges and struggles to replicate real-world visual complexity. Real-to-Sim methods~\cite{zhang2023real2sim, torne2024rialto} reconstruct specific real scenes in simulation, but are labor-intensive and tied to particular environments. TRANSIC~\cite{jiang2023transic} refines simulation-trained policies via online human corrections, while DrEureka~\cite{ma2024dreureka} uses LLMs to automate randomization range design. V-Dreamer sidesteps these limitations by synthesizing a large variety of physically grounded scenes through generative diffusion models, naturally providing open-vocabulary visual diversity without manual specification or real-world interaction, and enabling zero-shot sim-to-real transfer to unseen objects.